\subsection{Comparaison des capacités spécifiques massiques de quelques électrodes courantes des batteries \ce{Li}}
%from Culture science chimie

\setcounter{equation}{0}
\setcounter{Q}{0}

La capacité spécifique massique est la charge électrique pouvant être délivrée par gramme de matière active, c'est-à-dire par gramme du matériau hôte. Dans le système international, on l'exprime en \si{\coulomb\per\gram} mais en pratique, on l'utilise plutôt en \si{\milli\ampere\hour\per\gram}.

\Q \label{Q:11-01} Définir un ampère-heure [\si{\ampere\hour}] et établir le facteur de conversion avec le coulomb [\si{\coulomb}].
\solution{Un ampère-heure est une unité qui représente la quantité de charge délivrée par un dispositif débitant 1 ampère pendant 1 heure. Or, \SI{1}{\hour} équivaut à \SI{3600}{\second}, on peut donc en déduire que : 
\begin{align*}
\SI{1}{\ampere\hour} = \SI{3600}{\coulomb}
\end{align*}}

\Q \label{Q:11-02} On va considérer une réaction rédox modèle de la forme : 
\begin{align*}
\ce{M + z Li^+ + z e^- = MLi_z(s)}
\end{align*}
où $M$ représente le matériau hôte du lithium.

\'{E}tablir que la capacité spécifique massique s'exprime en \si{\milli\ampere\hour\per\gram} et peut être estimée par la relation :  
\begin{align*}
C_\text{Hote} = \frac{z F}{3.6 M_{Hote}}
\end{align*}

où $z$ est le nombre d'électrons échangé, $F$ la constante de Faraday et $M_{Hote}$ la masse molaire du matériau hôte.

\solution{
On va considérer que l'on échange $n_e$ moles électrons et de lithium avec un matériau hôte de masse $m_{hote}$, c'est-à-dire que le matériau peut accueillir au maximum $n_e$ moles de lithium dans sa structure ! Dressons un tableau d'avancement de la situation : 
\begin{center}\begin{tabular}{cccccccc}
\si{\mole} & \ce{M(s)} & \ce{ + } & \ce{z Li^+} & \ce{ + } & \ce{z e^-} & \ce{ = } & \ce{MLi_z(s)} \\
$t=0$ & $\frac{m_{hote}}{M_{hote}}$ & & $n_e$ & & $n_e$ & & 0\\
$t=\infty$ & $\frac{m_{hote}}{M_{hote}} - \xi_{\infty}$ & & $n_e - z\xi_{\infty}$ & & $n_e - z \xi_{\infty}$ & & $\xi_\infty$\\
\end{tabular}\end{center}
Lorsque l'échange d'électrons est terminé, tous les électrons échangeables ont été transférés, donc : 
\begin{align*}
\begin{cases}
\frac{m_{hote}}{M_{hote}} - \xi_{\infty} &= 0 \\
n_e - z\xi_{\infty} = 0
\end{cases} \implies \frac{m_{hote}}{M_{hote}} = \frac{n_e}{z} \implies n_e = \frac{z m_{hote}}{M_{hote}}
\end{align*}
Ainsi, on a une relation simple entre la structure hôte et le nombre d'électrons échangé. Or, on sait exprimer la quantité d'électrons échangés en quantité de charges échangées $Q$ : 
\begin{align*}
Q = F n_e = \frac{z F m_{hote}}{M_{hote}} \quad [\si{\coulomb}]
\end{align*}
La capacité spécifique massique de cette structure correspond au nombre d'électrons échangés par gramme de matériau, que l'on peut exprimer en \si{\coulomb\per\gram} puis en \si{\ampere\hour\per\gram} grâce à la Q.\ref{Q:11-01} :
\begin{align*}
C_{hote} &= \frac{Q}{m_{hote}} = \frac{z F }{M_{hote}} \quad [\si{\coulomb\per\gram}] \\
&= \frac{z F}{3.6 M_{hote}} \quad [\si{\milli\ampere\hour\per\gram}]
\end{align*}
}

\Q Dans le cas des matériaux lithiés suivants, donner la demi-équation rédox de réduction du cation lithium en lithium métallique. En déduire la capacité spécifique et compléter le tableau :

\begin{center}
\begin{tabular}{ccccccc}
\toprule
Matériau & \ce{Li(s)} & \ce{LiC6(s)} & \ce{Li15Si4(s)} & \ce{LiCoO2(s)} & \ce{LiFePO4(s)} & \ce{LiNi_{1/3}Co_{1/3}Mn_{1/3}O2} \\
\midrule
Masse molaire [\si{\gram\per\mole}] & 6.94 & 12.0 & 28.1 & 97.9 & 158 & 96.5 \\
\midrule
Capacité [\si{\milli\ampere\hour\per\gram}] & & & & & & \\
%Capacité [\si{\milli\ampere\hour\per\gram}] & 3960 & 372 & 3600 & 274 & 170 & 279 \\
\bottomrule
\end{tabular}
\end{center}

Commenter le résultat. \emph{Aide : Vous penserez à choisir une st{\oe}chiométrie compatible avec l'équation modèle étudiée à la Q. \ref{Q:11-02}.}

\solution{On peut écrire toutes les équations rédox en réduction, suivant le modèle de l'énoncé : 
\begin{align*}
& \ce{Li^+ + e^- = Li} & z = 1\\
& \ce{\frac{1}{6} Li^+ + \frac{1}{6} e^- + C = \frac{1}{6} LiC6} & z = 1/6\\
& \ce{\frac{15}{4} Li + \frac{15}{4} e^- + Si = \frac{1}{4} Li_{15}Si4} & z = 15/4\\
& \ce{Li^+ + e^- + CoO2 = LiCoO2} & z = 1 \\
& \ce{Li^+ + e^- + FePO4 = LiFePO4} & z = 1 \\
& \ce{Li^+ + e^- + Ni_{1/3}Co_{1/3}Mn_{1/3}O2 = LiNi_{1/3}Co_{1/3}Mn_{1/3}O2} & z=1
\end{align*}
On peut alors compléter le tableau proposé : 
\begin{center}
\begin{tabular}{ccccccc}
\toprule
Matériau & \ce{Li(s)} & \ce{LiC6(s)} & \ce{Li15Si4(s)} & \ce{LiCoO2(s)} & \ce{LiFePO4(s)} & \ce{LiNi_{1/3}Co_{1/3}Mn_{1/3}O2(s)} \\
\midrule
Masse molaire [\si{\gram\per\mole}] & 6.94 & 12.0 & 28.1 & 97.9 & 158 & 96.5 \\
\midrule
Capacité [\si{\milli\ampere\hour\per\gram}] & 3860 & 372 & 3580 & 274 & 170 & 279 \\
\bottomrule
\end{tabular}
\end{center}
On observe qu'en général, les matériaux qui accueillent le lithium n'atteignent pas \SI{10}{\percent} des capacités du lithium métallique. Cependant, une exception intéressante apparaît dans la liste : \ce{Li15Si4}. En effet, les matériaux au silicium sont aujourd'hui des candidats prometteurs pour leur excellent rapport Li/Si (=très nombreux sites d'insertion) qui permettent d'approcher des performances en capacité du lithium métallique ! \\
\textbf{Nota Bene :} La capacité spécifique ne dépend pas de la "nature de la structure cristallographique" du matériau.}

\paragraph*{Données}
\begin{itemize}
\item Constante de Faraday : $F = 96485 ~ \si{\coulomb\per\mole}$ 
\end{itemize}

\begin{comment}
\begin{itemize}
\item Masses molaires des éléments : 
\begin{center}
\begin{tabular}{cccccccc}
Lithium & Carbone & Oxygène & Silicium & Cobalt & Fer & Manganèse & Phosphore \\
6.94 & 12.0 & 16.0 & 28.1 & 59.0 & 55.9 & 55.0 & 31.0
\end{tabular}
\end{center}
\end{itemize}
\end{comment}